%%%%%
%
\section{Arctangente}

\itemstyle{Op\'erateur~:} \operfmt{Atan2}
\label{OP-Atan2}
\index{Atan2}

%%%%%OBJET%%%%
\itemstyle{Objet~:}
Calcule l'arctangente du rapport de deux variables, en fournissant
un r�sultat dans la gamme $]-\pi, \pi[$ (ou $]-180, 180[$).
Cela revient � calculer l'arctangente de y/x,
puis  utiliser le signe de x pour 
d�terminer le quadrant du r�sultat.

Le r�sultat est la mesure de l'angle entre l'axe
des X d'un rep�re othonorm� du plan et le segment
joignant l'origine de ce rep�re au point de coordonn�es (x,y),
compt� positivement dans le sens trigonom�trique direct.

%%%%%ARGUMENTS%%%%
\bigskip
%%%%%ENTREES%%%%
\begin{tabarg}
\multicolumn{3}{l}{\titretab {Entr\'ees}}\\
\hline
$Y$ & Vecteur, \texttt{float} & Premier op�rande (unit� quelconque)\\
$X$ & Vecteur, \texttt{float} & Second op�rande (m�me unit� que $Y$)\\
unit & Coeff., \texttt{String} & Unit�  souhait�e pour le r�sultat (``degre'' ou ``radian'')\\
\end{tabarg}

\medskip
%%%%%SORTIES%%%%
\begin{tabarg}
\multicolumn{3}{l}{\titretab {Sorties}}\\
\hline
$\alpha$ & Vecteur, \texttt{float} & Angle tel que $\tan \alpha = Y/X$ (unit� indiqu�e par l'argument unit)\\
\end{tabarg}


%%%%%FORMULE%%%%
\itemstyle{Formule}

%% Atan2
\[
\alpha = \arctan (Y/X)
\]



%%%%%COMMENTAIRE%%%%
\itemstyle{Commentaire~:}

Attention � l'ordre des arguments : $Y$ est le premier argument.

La valeur du r�sultat pour le cas $X = Y = 0$ n'est 
pas d�finie math�matiquement. Cet op�rateur suit la
convention de la biblioth�que math�matique du langage C, pour
qui atan2(0,0) = 0, afin d'�viter de cr�er de fausses
valeurs manquantes lors d'op�rations
de conversion de coordonn�es rectangulaire/polaires.


%%%%%AUTEUR%%%%
\itemstyle{Auteur~:}
CNRM/GMEI/TRAMM

%%%%%REFERENCE%%%%
%\itemstyle{R\'ef\'erences~:}
%INCONNU
